\histitem{局部紧交换群}

到20世纪30年代初，谐波分析与群表示之间的互动已经确立。接下来的十年中，二者相互推动，互动大为加强，分别促成对方迅速发展。

这种互动的一个极其肥沃的土壤是几乎周期函数的研究；H. 玻尔于1924年在实变量函数的情形中引入了它们 {[}9{]}。根据他的一个基本定理，这些函数正是 \(\mathbf{R}\) 上可由 \(x \mapsto e^{i\lambda x}\) 、\(\lambda \in \mathbf{R}\) 的线性组合一致逼近的有界连续函数（参见 II，第292页，习题54）。博赫纳 {[}6{]} 于1926年证明，\(f \in \mathcal{C}_b(\mathbf{R})\) 是几乎周期的，当且仅当其平移 \(x \mapsto f(x - a)\) 对于 \(a \in \mathbf{R}\) 在一致收敛拓扑下构成 \(\mathcal{C}_b(\mathbf{R})\) 的相对紧子集。博赫纳和冯·诺伊曼注意到，若 G 是拓扑群，这一刻画显然在 G 上（右侧或左侧）提供了几乎周期函数的概念。20世纪30年代初，在冯·诺伊曼 {[}52{]} 及彼得和外尔的方法推动下，这一理论有望迅速发展。

另一个富有成果的领域出现在 N. 维纳工作的边缘。他在1932年一篇著名论文 {[}88{]} 中把代数方法引入所谓塔伯尔问题的研究——也就是在已知某些加权平均行为的情况下，研究函数（或序列）沿滤子的渐近行为。他的方法受到 R. 施密特的思想启发。

维纳工作中的一个辅助结果（{[}88，引理 IIe{]}）尤其令人印象深刻并激发了诸多发展：若 \(f\) 是一个傅里叶级数绝对收敛且 \(f\) 不为零的周期函数，则 \(1/f\) 的傅里叶级数绝对收敛（参见 I，第38页，例）。早在1932年，R. 佩利和 N. 维纳 {[}89{]} 就勾勒出这一收敛结果、离散交换群的对偶群概念与几乎周期函数理论之间的联系。

在这一分析背景下，代数拓扑促使 L. 庞特里亚金为局部紧阿贝尔群的特征标研究提供决定性推动。他早在1932年就宣布，特征标群概念的意义在于把欧氏空间紧子集的同调群与其补集的同调群置于对偶关系中 {[}55{]}。在1934年发表的论文 {[}58{]} 中，他证明离散群的特征标群是紧的，随后在这一框架中建立对偶定理。虽然庞特里亚金使用了彼得和外尔的结果以及刚刚在一般情形建立存在性的哈尔测度，但他的主要目标是拓扑应用 {[}57{]}，并未明确关注谐波分析 \(^{(8)}\)：这些对偶定理是通过细致研究紧阿贝尔群的结构证明的。

庞特里亚金的工作立即引起极大兴趣。应冯·诺伊曼的建议，E. 范坎彭次年将对偶性推广到所有局部紧交换群 {[}37{]}。不久之后，他把新不变分析应用到几乎周期函数的研究，给出一个引人注目的结果 {[}38{]}；A. 外尔在同一时期独立发现了该结果 {[}79{]}：若 G 是局部紧阿贝尔群，可以得到一个紧群 \(\widetilde{G}\)：先给对偶群 \(\widehat{G}\) 赋予离散拓扑，再把由此得到的离散群的紧对偶群定义为 \(\widetilde{G}\)。群 G 随后典范地嵌入 \(\widetilde{G}\)（参见 II，第292页，习题54），对 \(\widetilde{G}\) 上连续函数应用彼得--外尔理论的简单结果，就能得到当时关于 G 上几乎周期函数的所有已知结果。

与此同时，在这一时期，经典分析中早已熟知的正定函数出现在 R 的酉表示研究中。J. 默瑟在1910年以前不久、结合积分方程理论研究了普遍正核 {[}45{]}；而其序列类比则在同一时期引出大量著作，尤其与矩问题有关（参见 IV，第359页，习题21）。与此同时，M. 马蒂亚斯于1923年已系统研究 R 上的正定函数 {[}43{]}；在博赫纳手中，它们成为傅里叶分析和概率论的基本工具之一 {[}7{]}。1933年，博赫纳 {[}8{]} 与里斯 {[}61{]} 独立注意到，R 的酉表示的每个对角矩阵系数都是 R 上的正定函数。这一观察使他们能够给出斯通定理的更简单证明，并将对一般理论的发展产生决定性影响（参见第7号）。

这一时期的结果由 A. 外尔整理在一本1937年前完成、1940年出版的著作中 {[}80{]}。在那里，他详细阐述彼得--外尔、庞特里亚金和范坎彭的结果。但在交换群情形，庞特里亚金和范坎彭强调群的结构和对偶定理，而外尔系统地发展了谐波分析，引入傅里叶变换并证明普朗谢雷尔公式与博赫纳定理（V，第455页，定理5）。他尤其强调卷积（称为“合成乘积”）的作用，以及正定函数的作用；后者如今定义在任意群上，并同表示的对角矩阵系数联系起来，至少在有限维表示情形如此。

因此在交换情形达到的这种一般性，稍晚些时候将为数论开辟新的视野。早在1936年，为向类域论引入代数方法，C. 谢瓦莱 {[}14{]} 引入了数域的 idèle 群；该群后来会作为 adèle 环的可逆元群出现（参见 AC，VII，第221--222页及 ÉHM，第143页）。1950年，J. 泰特 {[}75{]} 将证明，对 adèle 和 idèle 群的谐波分析使得可以容易地恢复 Hecke L-函数的函数方程，预示了把 adèle 群上的谐波分析同自守形式研究联系起来的非凡发展。

可以说，到20世纪30年代末，不变谐波分析的主要结果已经在紧群和局部紧交换群上建立。在交换情形，外尔的证明仍像庞特里亚金和范坎彭的证明一样，依赖于对群结构的精细了解，粗略说即 II 第2节第244页的分类结果。下一个十年借助盖尔范德学派的新方法将说明如何摆脱这一依赖（参见 H. 嘉当和 R. 戈德芒 {[}12{]}）。

